\chapter{摇滚：把电流变成心跳}
\label{ch:genre-rock-deep}

\classquote{审一以定和，比物以饰节，节奏合以成文。}{《礼记·乐记》}{先定一个主音来建立和谐，再用不同乐器装饰节拍，让节奏组合成章法。}

\section{什么是摇滚：先有骨架，再有电流}
\label{sec:rock-what}

\term{摇滚}{rock}不是“把音量推大”的同义词。它更像一种把身体和情绪同时推向前方的写法：一条能记住的 riff、一个清楚的反拍、一层有重量的失真，再加上一次真的释放。它可以很粗糙，也可以很精致；可以只有三个人，也可以有键盘、铜管和第二条吉他。真正让它成立的，是骨架先站住，电流才有地方通过。

四个抓手：

\begin{longtable}{p{0.16\textwidth}p{0.32\textwidth}p{0.38\textwidth}}
\toprule
抓手 & 听感 & 写作问题 \\
\midrule
\endfirsthead
\toprule
抓手 & 听感 & 写作问题 \\
\midrule
\endhead
Riff & 一段短音型反复出现，像歌的“身份证” & 去掉这条 riff，歌还认得出自己吗？ \\
反拍 & 军鼓落在 2、4 拍，身体被推着走 & 反拍够不够清楚？底鼓有没有把 1、3 拍踩稳？ \\
失真 & 吉他/贝斯的谐波变厚，情绪有“电流感” & 失真是在承担主 riff，还是只是把所有轨糊在一起？ \\
释放 & 主歌压住，副歌/段尾真的打开 & 副歌比主歌多了什么？如果只是更响，够不够？ \\
\bottomrule
\end{longtable}

摇滚最动人的时候，往往不是“最响”的时候，而是“压了很久、终于放开”的时候。安静的主歌、只留贝斯和鼓的桥段、Solo 前的两拍空，都是为了让后面的电流有对比。没有收，就没有放。

\section{从车库到体育场：摇滚的来路}
\label{sec:rock-history}

摇滚的谱系不是一条直线。它从布鲁斯、乡村、节奏布鲁斯和福音歌里长出来，又被每一代年轻人用自己的设备和场地改写。下面这张表只抓“编制和传播方式变了，写法就跟着变”的主线。

\begin{longtable}{p{0.16\textwidth}p{0.34\textwidth}p{0.42\textwidth}}
\toprule
时期 & 关键变化 & 对编配的影响 \\
\midrule
\endfirsthead
\toprule
时期 & 关键变化 & 对编配的影响 \\
\midrule
\endhead
1950 年代 & 摇滚乐（rock and roll）进入青少年市场；Chuck Berry、Elvis Presley、Little Richard 等把电吉他、钢琴和背拍推到前台 & 鼓的 2、4 拍成为身体坐标；电吉他从伴奏升为标志音色；歌曲短、钩子直接 \\
1960 年代 & 英伦入侵（The Beatles、The Rolling Stones）、Jimi Hendrix 的反馈与失真、The Who 的强力演出 & 吉他音色成为作曲材料；录音室从“记录演出”变成“制造声音”；专辑开始容纳更长的结构 \\
1970 年代 & Led Zeppelin、Black Sabbath 的重型 riff；Pink Floyd 的前卫结构；Ramones、Sex Pistols 的朋克反弹 & 重型失真、低速/中速 riff、长段落与三分钟朋克并存；舞台音量与音箱成为乐器的一部分 \\
1980 年代 & MTV、金属、另类摇滚、U2 与 R.E.M. 等; 鼓机与合成器进入摇滚 & 视觉与音色一起塑造乐队身份；吉他墙与大厅混响成为“大”的手段 \\
1990 年代 & Nirvana、Pearl Jam、Soundgarden 的 grunge；Radiohead 等另类摇滚 & “静—响”动态成为结构；粗糙音色与真实鼓声回到前台；歌词与姿态同样重要 \\
2000 年代 & The White Stripes、Arctic Monkeys、车库/独立摇滚回潮；数字发行 & 少编制、强 riff、复古录音质感；两三个人也能做出完整的低频与冲击 \\
2010 年代至今 & 流媒体、音乐节、曲风混搭；摇滚与电子、嘻哈、流行互相借用 & riff 可以很短，音色可以很“非吉他”；结构更碎片化，但现场释放仍然是核心 \\
\bottomrule
\end{longtable}

华语摇滚常被追溯到 1980 年代：崔健在 1986 年唱出《一无所有》，把电吉他、乐队和个体表达带进大众视野；1990 年代唐朝、黑豹等把重金属与硬摇滚做成产业；2000 年代以后，独立乐队、音乐节与互联网让摇滚进入更分化的场景。不同地区的工业条件不同，但问题一样：怎样用有限的编制，做出清楚的反拍、riff 和释放。

\section{配器：让每件乐器咬住同一副骨架}
\label{sec:rock-craft}

摇滚的编制可以很少，但不能“各弹各的”。把角色写进轨道名，比堆满音箱模拟更有效。

\subsection{典型编制与五层角色}
\label{sec:rock-instrumentation}

\begin{longtable}{p{0.16\textwidth}p{0.25\textwidth}p{0.25\textwidth}p{0.20\textwidth}}
\toprule
功能层 & 常见乐器/轨道 & 主要任务 & 常见错误 \\
\midrule
\endfirsthead
\toprule
功能层 & 常见乐器/轨道 & 主要任务 & 常见错误 \\
\midrule
\endhead
人声 & 主唱、和声、喊唱 & 讲故事、给姿态、承担副歌的记忆点 & 用过多修音/混响把摇滚人声做成流行 \\
节奏吉他 & 失真/半失真吉他、强和弦 & 给 riff、给和弦墙、和贝斯锁低频 & 从头到尾满扫，主唱没有呼吸 \\
主音吉他 & 失真吉他、滑棒/反馈 & 做短句呼应、Solo、段尾钩子 & Solo 只求快，没有动机和留白 \\
贝斯 & 电贝斯、过载/失真 & 锁底鼓、补根音、让低频有“咬合” & 只跟吉他同音，低频变成一堵墙 \\
鼓 & 原声鼓组、打击 & 定反拍、定段落地板、给 fill & 每段都加花，副歌没有“到”的感觉 \\
\bottomrule
\end{longtable}

一个实用的起点是六条轨：主唱、节奏吉他、主音吉他、贝斯、鼓、Master 总线。先把五件乐器的关系跑通，再加第二吉他、键盘或打击。轨数不是摇滚的勋章；“每件乐器都知道自己咬在哪一拍”才是。

\subsection{节奏：反拍要清楚，底鼓要稳}
\label{sec:rock-groove}

摇滚最常见的骨架是 4/4、90—140 BPM；朋克更快，重型/厄运更慢。无论速度如何，军鼓的 2、4 拍要清楚，底鼓通常落在 1、3 拍，或与贝斯的关键 riff 同落。

三条经验：

\begin{itemize}
\item \textbf{先锁底鼓与贝斯}：贝斯的根音落在底鼓上， riff 的强音可以让两者一起落；不要每一拍都让贝斯跟着吉他跑。
\item \textbf{反拍只留一个主角}：军鼓已经是反拍的主角，节奏吉他不要在每个 2、4 拍都加同样的重音，否则反拍会糊。
\item \textbf{fill 留给段尾}：鼓的 fill 放在主歌进副歌、Solo 进副歌、尾奏前；每四小节都来一次，听众会疲劳。
\end{itemize}

\subsection{吉他：一条 riff，三种工作}
\label{sec:rock-guitar}

电吉他在摇滚里通常做三件事，不要让它同时做三件：

\begin{enumerate}
\item \textbf{Riff}：短、低、重复，通常用\term{强和弦}{power-chord}（根音＋五度）或单音；它负责记忆点。
\item \textbf{Comping}：在主唱句尾给和弦，或在前奏/副歌做八分扫弦；它负责和声与推动。
\item \textbf{Lead}：短句呼应、Solo、反馈/噪音；它负责“说话”和释放。
\end{enumerate}

失真不是越多越好。主 riff 用一组失真，节奏吉他可以用稍干净的音色做对比，主音吉他在副歌让位。两条失真吉他从头到尾叠在一起，低频会糊，主唱会失去位置。

\subsection{贝斯与鼓：低频的咬合}
\label{sec:rock-rhythm-section}

贝斯不是“低八度的吉他”。它要和底鼓一起把低频的地基打出来：

\begin{itemize}
\item 根音落在强拍，经过音只用在换和弦前一拍；
\item 如果吉他 riff 在高音区，贝斯可以走低音区的对应动机，不要完全同音；
\item 失真贝斯可以增加中频的“咬”，但不要把低频推到无法与底鼓分离；
\item 鼓的底鼓要短、实，军鼓要有反拍的“啪”；混响不要大到让反拍变软。
\end{itemize}

\subsection{人声与结构：压住，再放开}
\label{sec:rock-vocal-form}

摇滚人声不需要一直喊。主歌可以用较低的 chest voice、近讲、少混响；副歌再把力度和空间打开。关键是让“放开”有对比。

常见结构：

\begin{itemize}
\item 前奏 riff → 主歌（riff＋人声）→ 副歌（吉他加厚、鼓更重）→ 主歌 → 副歌 → 桥段（换和声/降密度）→ Solo → 副歌 → 尾奏。
\item 也可以把 riff 当段落标记：前奏、副歌入口、尾奏都用同一条 riff，听众会知道“回来了”。
\item Solo 不是“多弹音符”，而是把主歌的动机放大；Solo 段可以只留主音吉他、贝斯和鼓，给耳朵腾位置。
\end{itemize}

\section{四首歌的配器拆解}
\label{sec:rock-cases}

以下只做编配分析，不引歌词、不附旋律谱例。作品信息以专辑发行信息与官方视频页为据；链接见章末扩展阅读。

\subsection{Queen《Bohemian Rhapsody》：把一首歌写成几幕}
\label{sec:rock-case-bohemian}

Queen 的《Bohemian Rhapsody》1975 年发行，收录于《A Night at the Opera》。它不像常规摇滚那样靠一条 riff 贯穿，而是把钢琴、合唱、吉他、歌剧式段落拼成一首“多幕剧”。

\begin{itemize}
\item \textbf{结构}：前奏的钢琴合唱、叙事段、吉他段、歌剧段、硬摇滚段、尾奏，每一段都有不同的密度和音色；听众被带着走，而不是等一个副歌。
\item \textbf{配器}：钢琴承担和声骨架，人声叠唱制造“合唱团”，吉他在硬摇滚段才真正压上来；每件乐器在各自段落里当主角。
\item \textbf{编配逻辑}：用“段落对比”代替“循环叠加”。这和常规流行/摇滚的写法不同，但说明了一个原则：riff 不是唯一的记忆点，段落本身也可以成为记忆点。
\end{itemize}

可借鉴：如果你的歌有几段不同的情绪，不要强行用同一个副歌循环；可以设计一个“回到主 riff/主句”的锚，让听众在变化里有路可回。

\subsection{Nirvana《Smells Like Teen Spirit》：静—响是结构}
\label{sec:rock-case-teen-spirit}

Nirvana 的《Smells Like Teen Spirit》1991 年发行，收录于《Nevermind》。它最值得学的是“安静的主歌与爆发的副歌”之间的动态设计。

\begin{itemize}
\item \textbf{吉他}：主 riff 用低音区、强和弦与失真；主歌压住，副歌把同一动机放大，而不是换一条全新的素材。
\item \textbf{鼓}：军鼓的 2、4 拍与底鼓形成清楚的背拍；副歌的 crash 与更密的鼓让“放开”有物理重量。
\item \textbf{人声}：主歌近、低、少混响；副歌放开，但不靠过度修音；歌词的强词落在反拍/强拍上。
\item \textbf{编配逻辑}：同一组材料，靠密度、失真度和人声力度做“静—响”；这让副歌的爆发不需要额外堆很多轨。
\end{itemize}

可借鉴：先写一个能在低密度下成立的 riff，再决定副歌加什么。若主歌已经把所有东西都用完，副歌只能更响，不能更“开”。

\subsection{The White Stripes《Seven Nation Army》：一条 riff，两个人和一座体育场}
\label{sec:rock-case-seven-nation}

The White Stripes 的《Seven Nation Army》2003 年发行，收录于《Elephant》。它用一条低音区 riff、极简的鼓和两个人，做出了体育场级别的记忆点。

\begin{itemize}
\item \textbf{Riff}：那条像贝斯一样的吉他 riff 是全曲的主钩子；它低、短、重复，现场几万人可以跟着唱。
\item \textbf{编制}：只有吉他、鼓和人声，没有贝斯；riff 自己承担低频，鼓用极简的 kick/snare 把反拍钉住。
\item \textbf{空间}：主歌留空，副歌加人声/吉他叠层；不是靠更多乐器，而是靠“同一 riff 在不同密度下重复”。
\item \textbf{编配逻辑}：少编制不是限制，而是把每个声音的职责放大。riff 清楚，鼓清楚，人声清楚，歌就成立。
\end{itemize}

可借鉴：如果你只有两三条轨，不要急着用 Pad 填满；先问“这条 riff 能不能同时承担记忆点和低频”。能，就把它做清楚。

\subsection{Arctic Monkeys《Do I Wanna Know?》：riff 也可以很“冷”}
\label{sec:rock-case-do-i-wanna-know}

Arctic Monkeys 的《Do I Wanna Know?》2013 年发行，收录于《AM》。它把摇滚 riff 做得偏冷、偏慢、偏夜晚，而不是靠高增益失真。

\begin{itemize}
\item \textbf{Riff}：低音区的吉他动机反复出现，带一点滑音和空弦；它的“酷”来自音色和留白，不是来自速度。
\item \textbf{律动}：鼓的背拍清楚但不急躁，贝斯跟着 riff 的低音走；整体像半速的摆动，给主唱留出空间。
\item \textbf{人声}：主唱偏近、偏低，副歌用假声/和声打开；人声不是喊，而是把情绪压在里面。
\item \textbf{编配逻辑}：riff 负责记忆点，鼓/贝斯负责摆动，人声负责情绪；三者各做一件事，所以听起来空，却不薄。
\end{itemize}

可借鉴：失真不是摇滚的唯一答案。把 riff 的音色做窄、把低频留出来、把鼓的反拍做清楚，也可以很摇滚。

\section{夜半骨架：八小节摇滚}
\label{sec:rock-yeban}

下面是一套可以在夜半重建的八小节骨架，重点是 riff、反拍和“静—响”。

\begin{longtable}{p{0.10\textwidth}p{0.22\textwidth}p{0.58\textwidth}}
\toprule
小节 & 段落 & 编配动作 \\
\midrule
\endfirsthead
\toprule
小节 & 段落 & 编配动作 \\
\midrule
\endhead
1—4 & 主歌 & 节奏吉他：低音区 riff，强和弦/单音；贝斯：根音锁底鼓；鼓：kick 1/3、snare 2/4、闭镲八分；主唱：中低音区，句尾留白；主音吉他：不进入 \\
5—8 & 副歌 & 节奏吉他：riff 加八分扫弦/双轨；贝斯：根音＋经过音；鼓：crash 与更密的 fill；主唱：放开，加一层和声；主音吉他：短句呼应，不盖主唱 \\
\bottomrule
\end{longtable}

操作顺序：

\begin{enumerate}
\item 新建 100 BPM、4/4 工程；建主唱、节奏吉他、主音吉他、贝斯、鼓五条轨，Master 保留为总线。
\item 先写节奏吉他的两小节 riff（例如 E5–G5–A5 的强和弦循环），只保留低音区；用 F7 钢琴卷帘修好每拍的落点。
\item 写鼓：底鼓 1/3、军鼓 2/4、闭镲八分；第 4 小节末加一个两拍 fill，送进副歌。
\item 写贝斯：根音跟底鼓，riff 的强音与吉他同落；不要在每一拍都跟着吉他跑。
\item 写主唱：主歌只用中低音区，句尾留半拍到一拍；第 5 小节副歌放开，但不要把所有句子都推到最高。
\item 第 5 小节加主音吉他的短句、crash 和一层和声；只加一层，确保主唱仍然清楚。
\item 用 F8 看电平，Solo 每条轨检查低频：贝斯与底鼓要咬住，节奏吉他与主音吉他要错开中频。
\item 用 F3 打开设备堆叠对照音色；如果失真把主唱盖住，先减失真/降低节奏吉他，不要只推主唱。
\item 保存工程，导出八小节试听；问自己“主歌压住了吗？副歌真的打开了吗？”
\end{enumerate}

\section{练习：三条路，不同取舍}
\label{sec:rock-practice}

\begin{exercise}
\textbf{两小节 riff，三种写法。} 用同一组 E–G–A 做三个版本：A 单音低音 riff；B 强和弦（根音＋五度）；C 强和弦＋手掌闷音。比较三版的“记忆点、低频重量、与主唱的让位”。A 清楚但可能薄；B 有重量但容易糊；C 有咬合但需要控制闷音长度。
\end{exercise}

\begin{exercise}
\textbf{静—响，三种抬升。} 给同一段主歌接三个副歌：A 只把鼓加满；B 只把吉他失真打开；C 先撤掉主音吉他/和声，再在副歌第一拍全部放回。逐个播放，记录哪一种“放开”最有对比；三种都试，再选适合你歌的那一种。
\end{exercise}

\begin{exercise}
\textbf{ riff 与鼓的咬合。} 把同一条 riff 分别对在底鼓的 1/3 拍、反拍的 2/4 拍、以及弱拍的十六分位置上；每次只改鼓，不改 riff。听哪一种最“推”，哪一种最“拖”。把结果写成一句话，贴在工程备注里。
\end{exercise}

\section{章末扩展阅读}
\label{sec:rock-reading}

\begin{sloppypar}
\begin{itemize}
\item \textbf{AllMusic, “Pop/Rock”}：\url{https://www.allmusic.com/genre/pop-rock-ma0000002613}。用流派条目补齐摇滚的谱系、子类型与代表艺人，适合先建立坐标系。
\item \textbf{Queen, “Bohemian Rhapsody” (Official Video Remastered)}：\url{https://www.youtube.com/watch?v=fJ9rUzIMcZQ}。看一首歌如何用多段落、多音色完成“几幕剧”式的编配。
\item \textbf{Nirvana, “Smells Like Teen Spirit” (Official Music Video)}：\url{https://www.youtube.com/watch?v=hTWKbfoikeg}。对照安静主歌与爆发副歌之间的密度、失真和鼓的层次。
\item \textbf{The White Stripes, “Seven Nation Army” (Official Music Video)}：\url{https://www.youtube.com/watch?v=0J2QdDbelmY}。听两个人如何用一条 riff 和极简鼓撑起完整的低频与记忆点。
\end{itemize}
\end{sloppypar}

以上链接在 2026 年 11 月 27 日按官方页面与搜索结果逐条核对标题与作者；若平台地区限制导致无法打开，请以条目名在对应平台的官方账号内检索。

\section{本章要点}
\label{sec:rock-takeaways}

\begin{itemize}
\item 摇滚不是音量，而是 riff、反拍、失真与释放的组合。
\item 摇滚的演化跟着传播媒介和现场变化走：电台、电视、MTV、流媒体、音乐节，每一代都改变“什么声音能代表乐队”。
\item 编制先写功能，再写乐器；节奏吉他、主音吉他、贝斯、鼓各做一件事。
\item 反拍要清楚，底鼓要稳；fill 留给段尾，不要每四小节都来一次。
\item 失真不是越多越好；让主 riff 有重量，让人声有位置。
\item 静—响是最可靠的摇滚结构之一：主歌压住，副歌才有得放。
\item 夜半骨架按“riff—鼓/贝斯—主唱—主音/色彩—电平检查”的顺序落地；能删、能解释，才是好编配。
\end{itemize}

\begin{helpmapping}
\begin{longtable}{p{0.45\textwidth}p{0.45\textwidth}}
\toprule
小节 & 帮助锚点 \\
\midrule
什么是摇滚 & \texttt{arrangement-genre-rock-what} \\
从车库到体育场 & \texttt{arrangement-genre-rock-history} \\
配器：让每件乐器咬住同一副骨架 & \texttt{arrangement-genre-rock-craft} \\
四首歌的配器拆解 & \texttt{arrangement-genre-rock-cases} \\
夜半骨架 & \texttt{arrangement-genre-rock-yeban} \\
练习 & \texttt{arrangement-genre-rock-practice} \\
章末扩展阅读 & \texttt{arrangement-genre-rock-reading} \\
\bottomrule
\end{longtable}
\end{helpmapping}

\begin{figure}[htbp]\centering
\begin{tikzpicture}[font=\scriptsize,x=1cm,y=1cm,>=Stealth]
\yebangenreddraw{摇滚：让失真与律动同时站住}{60--200 Hz 吉他低频}{200--2 kHz 失真中频}{2--8 kHz 镲片与瞬态}
\end{tikzpicture}
\caption{摇滚频段分工、鼓组网格与段落能量曲线}
\end{figure}
